\documentclass[11pt]{article}
\usepackage{amsmath,amssymb,amsthm}
\usepackage{hyperref}

\newtheorem{theorem}{Theorem}
\newtheorem{proposition}[theorem]{Proposition}
\newtheorem{lemma}[theorem]{Lemma}
\newtheorem{conjecture}[theorem]{Conjecture}
\newtheorem{definition}[theorem]{Definition}
\newtheorem{remark}[theorem]{Remark}
\newtheorem{example}[theorem]{Example}
\newtheorem{observation}[theorem]{Observation}

\newcommand{\N}{\mathbb{N}}
\newcommand{\Z}{\mathbb{Z}}
\newcommand{\Q}{\mathbb{Q}}
\newcommand{\CPP}{\operatorname{CPP}}
\newcommand{\GK}{\operatorname{GK}}

% q-series notation
\newcommand{\qpoch}[2]{(#1;#2)_\infty}
\newcommand{\qpochn}[3]{(#1;#2)_{#3}}
\newcommand{\qbin}[2]{\genfrac{[}{]}{0pt}{}{#1}{#2}}

% QSym notation
\newcommand{\QSym}{\operatorname{QSym}}
\newcommand{\Sym}{\operatorname{Sym}}

\title{Quasi-symmetric refinement of the cylindric partition\\
generating function}
\author{Robin Langer}
\date{\today}

\begin{document}
\maketitle

\begin{abstract}
We show that the generating function for rank-3 cylindric plane
partitions, refined by value frequencies, is a quasi-symmetric
function. Specifically, for a profile $c = (c_0, c_1, c_2)$,
the power series
$\widetilde{\GK}_c(\mathbf{z}) = \sum_{\pi \in \CPP(c)}
z_1^{\#_1(\pi)} z_2^{\#_2(\pi)} \cdots$
(where $\#_k(\pi)$ counts the number of parts equal to~$k$ in~$\pi$)
lies in the ring $\QSym$ of quasi-symmetric functions.
This is verified computationally for all profiles with
$d = |c| \leq 5$. Under the principal specialization
$z_k \mapsto q^k$, this recovers
$\GK_c(1, q) = \sum_n g_n(q)$.
Moreover, the expansion in Gessel's fundamental basis
$\{L_\alpha\}$ has nonneg coefficients ($L$-positivity),
a strictly stronger property than $M$-positivity.
However, the quasi-symmetry applies only to the
\emph{unconstrained} generating function (all CPPs, no
bound on the maximum entry). Since Warnaar's
$Q_{n,c}(q)$ polynomials are extracted from the
$\max = n$ piece, the QSym structure does not directly
bear on the positivity conjecture.
\end{abstract}

%% ========================================================================
\section{Setup}
\label{sec:setup}
%% ========================================================================

A \emph{rank-3 cylindric plane partition} with profile
$c = (c_0, c_1, c_2)$ is a triple of partitions
$\pi = (\nu^{(1)}, \nu^{(2)}, \nu^{(3)})$ satisfying the
shifted dominance (interlacing) conditions:
\begin{equation}\label{eq:interlacing}
  \nu^{(1)}_j \geq \nu^{(2)}_{j+c_1}, \qquad
  \nu^{(2)}_j \geq \nu^{(3)}_{j+c_2}, \qquad
  \nu^{(3)}_j \geq \nu^{(1)}_{j+c_0},
\end{equation}
for all $j \geq 0$, where $\nu^{(i)}_j = 0$ for
$j \geq \ell(\nu^{(i)})$.

The two-variable generating function is
\[
  \GK_c(z, q)
  = \sum_{\pi \in \CPP(c)} z^{n(\pi)}\, q^{|\pi|}
  = \sum_{n \geq 0} g_n(q)\, z^n,
\]
where $n(\pi) = \max\{\nu^{(i)}_j : i, j\}$ is the maximum entry
(equivalently, the number of paths in the NILP model) and
$|\pi| = |\nu^{(1)}| + |\nu^{(2)}| + |\nu^{(3)}|$ is the total weight.

\medskip

Warnaar's Conjecture~2.7 asserts that the polynomials
\[
  Q_{n,c}(q)
  = \qpochn{q}{q}{n} \cdot [z^n]\bigl(\qpoch{zq}{q} \cdot \GK_c(z,q)\bigr)
\]
have nonneg integer coefficients for all $n \geq 0$ and all profiles~$c$
with $d = |c| \not\equiv 0 \pmod{3}$.

%% ========================================================================
\section{The frequency refinement}
\label{sec:frequency}
%% ========================================================================

For a cylindric partition $\pi = (\nu^{(1)}, \nu^{(2)}, \nu^{(3)})$,
define the \emph{value frequency} of~$k$ as
\[
  \#_k(\pi) = \bigl|\{(i,j) : \nu^{(i)}_j = k\}\bigr|,
\]
the total number of parts equal to~$k$ across all three partitions.
The weight decomposes as $|\pi| = \sum_{k \geq 1} k \cdot \#_k(\pi)$.

\begin{definition}[Frequency generating function]
For a profile $c$, define
\begin{equation}\label{eq:GKtilde}
  \widetilde{\GK}_c(\mathbf{z})
  = \sum_{\pi \in \CPP(c)}
    z_1^{\#_1(\pi)}\, z_2^{\#_2(\pi)} \cdots z_n^{\#_n(\pi)} \cdots
  \;\in\; \Z[[z_1, z_2, \ldots]],
\end{equation}
where the sum is over all cylindric partitions with profile~$c$
(of any maximum entry).
\end{definition}

\begin{remark}
Under the specialization $z_k \mapsto q^k$, we recover
\[
  \widetilde{\GK}_c(q, q^2, q^3, \ldots)
  = \sum_{\pi} q^{|\pi|}
  = \GK_c(1, q).
\]
The maximum entry $n(\pi) = \max\{k : \#_k(\pi) > 0\}$ is
recoverable from the frequency data but is not a linear
statistic, so $z^{n(\pi)}$ is not a monomial specialization
of the~$z_k$.
\end{remark}

%% ========================================================================
\section{Quasi-symmetry}
\label{sec:qsym}
%% ========================================================================

Recall that a formal power series $f \in \Z[[z_1, z_2, \ldots]]$ is
\emph{quasi-symmetric} if for every composition
$\alpha = (\alpha_1, \ldots, \alpha_k)$ with $\alpha_i > 0$,
the coefficient of $z_{i_1}^{\alpha_1} \cdots z_{i_k}^{\alpha_k}$
is the same for all choices of indices $1 \leq i_1 < i_2 < \cdots < i_k$.

The ring $\QSym$ of quasi-symmetric functions sits between
the ring $\Sym = \Lambda$ of symmetric functions and the full power
series ring:
\[
  \Sym \;\subset\; \QSym \;\subset\; \Z[[z_1, z_2, \ldots]].
\]
The standard monomial basis of $\QSym$ is
$\{M_\alpha\}_{\alpha}$, indexed by compositions, where
\[
  M_\alpha
  = \sum_{i_1 < i_2 < \cdots < i_k}
    z_{i_1}^{\alpha_1} z_{i_2}^{\alpha_2} \cdots z_{i_k}^{\alpha_k}.
\]

\begin{observation}[Computational]\label{obs:qsym}
For every rank-3 profile $c$ with $d = |c| \leq 5$ (55 profiles total),
$\widetilde{\GK}_c(\mathbf{z})$ is quasi-symmetric. That is,
\begin{equation}\label{eq:qsym-expansion}
  \widetilde{\GK}_c(\mathbf{z})
  = \sum_{\alpha} a_\alpha(c)\, M_\alpha,
\end{equation}
with all coefficients $a_\alpha(c) \in \Z_{\geq 0}$.
\end{observation}

The nonnegativity of the coefficients $a_\alpha(c)$ is automatic
(they count cylindric partitions), so the content of the observation
is the quasi-symmetry itself: the count depends only on the
\emph{composition} of value frequencies, not on which values carry
those frequencies.

\begin{remark}
$\widetilde{\GK}_c$ is \emph{not} symmetric. For example,
with $c = (1,1,0)$, the coefficient of $M_{(2,1)}$ is~4
while the coefficient of $M_{(1,2)}$ is~3. Since
$(2,1)$ and $(1,2)$ are different compositions (related by
rearrangement, hence giving the same partition), symmetry
would require these to be equal.
\end{remark}

%% ========================================================================
\section{Monomial QSym coefficients for small profiles}
\label{sec:data}
%% ========================================================================

\subsection{Profile $c = (1,1,0)$, $d = 2$}

The QSym coefficients $a_\alpha = a_\alpha(1,1,0)$ display a
remarkable stabilization: once all parts of the composition~$\alpha$
are $\geq 2$, the coefficient depends only on the length of~$\alpha$.

\begin{center}
\begin{tabular}{c|l}
  \textbf{Composition $\alpha$} & \textbf{$a_\alpha$} \\
  \hline
  $()$ & $1$ \\
  $(k)$, any $k \geq 1$ & $2$ \\
  $(1, k)$, any $k \geq 1$ & $3$ \\
  $(j, k)$, $j \geq 2$ & $4$ \\
  $(1, 1, k)$, any $k \geq 1$ & $5$ \\
  $(1, j, k)$, $j \geq 2$ & $6$ \\
  $(j, 1, k)$, $j \geq 2$ & $6$ \\
  $(j, k, 1)$, $j,k \geq 2$ & $8$ \\
  $(j, k, l)$, $j,k \geq 2$, $l \geq 2$ & $8$ \\
  $(1, 1, 1, k)$ & $8$ \\
  $(1, 1, j, k)$, $j \geq 2$ & $10$ \\
  $(1, j, 1, k)$, $j \geq 2$ & $9$ \\
  $(j, 1, 1, k)$, $j \geq 2$ & $10$ \\
  $(1, j, k, l)$, $j,k \geq 2$ & $12$ \\
  $(j, 1, k, l)$, $j \geq 2$, $k \geq 2$ & $12$ \\
  $(j, k, 1, l)$, $j,k \geq 2$ & $12$ \\
  $(j, k, l, m)$, $j,k,l \geq 2$ & $16$
\end{tabular}
\end{center}

The stabilized values (all parts $\geq 2$) by length $\ell$ are:
$1, 2, 4, 8, 16, \ldots = 2^\ell$. This suggests a product
structure in the stable regime.

\subsection{Profile $c = (2,0,0)$, $d = 2$ (second Rogers--Ramanujan)}

The stabilized coefficients again follow powers of 2:
length $\ell$ gives $a_\alpha = 2^\ell$ when all parts of $\alpha$
are sufficiently large.

\subsection{Profile $c = (2,1,1)$, $d = 4$}

The coefficients are richer: $a_{(k)} = 5$ for $k$ large;
$a_{(j,k)} = 25$ for $j,k$ large;
$a_{(j,k,l)} = 125$ for $j,k,l$ large.
The stabilized values appear to be $5^\ell$, where $5$ is the
number of $A_2$ cylindric partitions of the empty profile (i.e.,
ordinary plane partitions of size~$\leq$~something related to~$d$).

%% ========================================================================
\section{Interpretation and structure}
\label{sec:interpretation}
%% ========================================================================

\subsection{Why quasi-symmetric?}

The quasi-symmetry has a natural explanation. Consider a CPP
$\pi$ with frequency composition $\alpha = (a_1, \ldots, a_k)$
at values $v_1 < v_2 < \cdots < v_k$. The interlacing
conditions~\eqref{eq:interlacing} compare entries at the same
positions: $\nu^{(i)}_j \geq \nu^{(i')}_j$ (with shifts).
Whether a particular entry equals $v_s$ or $v_{s'}$ matters for
the \emph{weight} but the combinatorial constraints depend only on
the \emph{relative order} of entries, not on their absolute values.

More precisely: the set of CPPs with value set $\{v_1, \ldots, v_k\}$
and frequency vector $(a_1, \ldots, a_k)$ is in weight-preserving
bijection with those using value set $\{w_1, \ldots, w_k\}$
(any other increasing sequence) with the same frequency vector,
via the substitution $v_s \mapsto w_s$ applied to every entry.
This bijection preserves the interlacing conditions because they
only depend on the ordering of entries.

This is precisely the defining property of quasi-symmetry: the
coefficient of $z_{v_1}^{a_1} \cdots z_{v_k}^{a_k}$ depends only
on $\alpha = (a_1, \ldots, a_k)$, not on $v_1 < \cdots < v_k$.

The argument is an instance of Stanley's $P$-partition
theory~\cite{stanley-ordered, gessel-multipartite}:
whenever a set of fillings is constrained by conditions
that only compare the \emph{relative order} of values
(weak and strict inequalities), the generating function
tracking value frequencies is quasi-symmetric.
The cylindric interlacing conditions~\eqref{eq:interlacing}
are precisely of this form.

Note that the individual
$\tilde{g}_n = \sum_{\max(\pi) = n} z_1^{\#_1} \cdots$
(fixed max entry~$= n$) is \emph{not} quasi-symmetric:
the coefficient of $z_n^1$ (a single entry equal to~$n$)
is positive, while $z_1^1$ (a single entry equal to~1, but
$\max = n$) is zero for $n > 1$.
Any constraint on the maximum entry, whether $\max = n$
or $\max \leq m$, breaks quasi-symmetry by privileging
certain values over others. Only the unconstrained
sum~\eqref{eq:GKtilde} is quasi-symmetric.

\begin{remark}
The function is not symmetric because swapping the
\emph{multiplicities} (e.g., having $a$ copies of value~1 and $b$
copies of value~2, versus $b$ copies of value~1 and $a$ copies of
value~2) changes the interlacing constraints: value~1 appears in
inner positions of the interlacing, while value~2 must dominate
value~1 entries. The relative \emph{positions} of the values in the
interlacing chain distinguish compositions from partitions.
\end{remark}

\subsection{Disconnect from the positivity conjecture}

The frequency generating function lives in $\QSym$, which is a
$\Sym$-module (in fact a $\Sym$-algebra). Natural bases of $\QSym$
include the monomial basis $\{M_\alpha\}$ (in which our
expansion~\eqref{eq:qsym-expansion} has manifestly nonneg
coefficients), the fundamental basis $\{L_\alpha\}$ of Gessel
(see Section~\ref{sec:L-positivity}), and the dual basis to
noncommutative symmetric functions.

However, this QSym structure does \emph{not} help with
Warnaar's Conjecture~2.7. The conjecture concerns the
polynomials $Q_{n,c}(q)$, which are extracted from
the coefficient of~$z^n$ in $\qpoch{zq}{q} \cdot \GK_c(z,q)$.
This extraction involves:
\begin{enumerate}
  \item The max-entry statistic $n(\pi) = \max\{\nu^{(i)}_j\}$,
    which is \emph{not} a monomial specialization of the~$z_k$
    (cf.\ Remark after~\eqref{eq:GKtilde});
  \item An alternating-sign sum (the $(zq;q)_\infty$ factor),
    whose cancellation crosses between different frequency vectors
    (see Section~\ref{sec:Qn-lift}).
\end{enumerate}
The quasi-symmetry of $\widetilde{\GK}_c$ is a property of the
\emph{aggregate} over all max entries, but $Q_{n,c}$ lives at
a \emph{fixed} max entry. The QSym framework cannot see this
decomposition.

In particular, the specialization $z_k \mapsto q^k$ sends
$\widetilde{\GK}_c$ to $\GK_c(1,q) = \sum_n g_n(q)$, which is
manifestly nonneg since each $g_n(q)$ counts objects. This tells
us nothing about the positivity of $Q_{n,c}(q)$, which involves
nontrivial cancellation among the~$g_n$.

\subsection{Partial quasi-symmetry of $\tilde{g}_n$}
\label{sec:partial-qsym}

Although $\tilde{g}_n$ (CPPs with $\max = n$) is not
quasi-symmetric in the variables $z_1, z_2, \ldots, z_n$
(since the value~$n$ is distinguished), it retains a
\emph{partial} quasi-symmetry: the value~$n$ is special,
but the remaining values $\{1, \ldots, n-1\}$ can be
freely relabelled.

More precisely, write
\[
  \tilde{g}_n
  = \sum_{m \geq 1} z_n^m \cdot f_{n,m}(z_1, \ldots, z_{n-1}),
\]
where $f_{n,m}$ collects those CPPs with exactly $m$ parts
equal to~$n$.

\begin{observation}[Computational]\label{obs:partial-qsym}
For every rank-3 profile $c$ with $d = |c| \leq 5$ and
$n = 1, 2, 3$: each slice $f_{n,m}$ is a quasi-symmetric
polynomial in~$z_1, \ldots, z_{n-1}$. That is,
\[
  \tilde{g}_n \;\in\; z_n \cdot \QSym(z_1, \ldots, z_{n-1})[z_n].
\]
\end{observation}

The argument is the same relabelling bijection as for the
full~$\widetilde{\GK}_c$, but restricted to non-$n$ values.
Given a CPP $\pi$ with $\max(\pi) = n$ and non-$n$ value
set $\{v_1 < \cdots < v_k\} \subset \{1, \ldots, n-1\}$,
replacing each $v_i$ by any other increasing sequence
$\{w_1 < \cdots < w_k\} \subset \{1, \ldots, n-1\}$
preserves the interlacing conditions: the non-$n$ entries
interact with each other only through their relative order,
and every non-$n$ entry~$v$ satisfies $v < n$, so the
inequalities involving $n$-entries are automatically preserved.

\begin{remark}[Failure of slice $L$-positivity]
The individual slices $f_{n,m}$ are $M$-positive
(the coefficients count CPPs) but are \emph{not} in general
$L$-positive in $\QSym(z_1, \ldots, z_{n-1})$.
Already for $c = (1,0,0)$, $n = 3$, $m = 1$
(two variables $z_1, z_2$), the $M$-coefficients
$a_{(k)} = 2$ and $a_{(j,k)} = 1$ yield
$L_{(a,b)} = M_{(a,b)} - M_{(a+b)} = 1 - 2 = -1$
for every two-part composition. This holds for all
55 tested profiles.
\end{remark}

\begin{remark}[Variable-count mismatch]
The E--H factorization $Q_n = \qpochn{q}{q}{n} \cdot g_n$
couples two objects that live in different numbers of
variables: the factor $\qpochn{q}{q}{n} = \prod_{k=1}^{n}(1-q^k)$
lifts to the elementary symmetric functions $e_1, e_2, \ldots$
in \emph{infinitely many} variables, while $\tilde{g}_n$
is a polynomial in only $z_1, \ldots, z_n$. This mismatch
is a structural obstacle to any multivariate lift of
$Q_{n,c}$: the cancellation producing $Q_n$ from $g_n$
necessarily involves objects living in incompatible
numbers of variables.
\end{remark}

\subsection{Stabilization and product structure}

For $c = (1,1,0)$, the stabilized QSym coefficients (all parts
of $\alpha$ sufficiently large) satisfy
$a_\alpha^{\mathrm{stab}} = 2^{\ell(\alpha)}$,
suggesting a product formula
\[
  \widetilde{\GK}_c^{\mathrm{stab}}
  = \prod_{k \geq 1} (1 + z_k)^?
\]
in some asymptotic sense. More generally, for $c = (2,1,1)$ with
$d=4$, the stable value is $5^{\ell(\alpha)}$, consistent with
the number of maximal cylindric configurations per value.
Understanding this product structure is a direction for future work.

%% ========================================================================
\section{$L$-positivity}
\label{sec:L-positivity}
%% ========================================================================

The \emph{fundamental quasi-symmetric functions} $L_\alpha$
(also denoted $F_\alpha$ in the literature) form a second
natural basis of $\QSym$, related to the monomial basis by
\[
  L_\alpha = \sum_{\beta \preceq \alpha} M_\beta,
\]
where $\beta \preceq \alpha$ means $\beta$ is a refinement
of~$\alpha$ (obtained by splitting parts of~$\alpha$).
Equivalently,
\[
  L_{(\alpha_1, \ldots, \alpha_k)}
  = \sum_{\substack{i_1 \leq i_2 \leq \cdots \leq i_n \\
    i_j < i_{j+1} \text{ if } j \in D(\alpha)}}
  z_{i_1} z_{i_2} \cdots z_{i_n},
\]
where $n = |\alpha|$ and $D(\alpha) = \{\alpha_1, \alpha_1+\alpha_2,
\ldots, \alpha_1+\cdots+\alpha_{k-1}\}$ is the descent set of~$\alpha$.

The fundamental basis plays the role in $\QSym$ that Schur functions
play in $\Sym$: $L$-positivity is the ``right'' notion of positivity
in $\QSym$. Under the natural map $\QSym \to \Sym$ that sends
$L_\alpha \mapsto s_{\operatorname{sort}(\alpha)}$,
$L$-positivity implies Schur-positivity of the symmetrization.

\begin{observation}[Computational]\label{obs:L-pos}
For every rank-3 profile $c$ with $d = |c| \leq 5$ (55 profiles total),
$\widetilde{\GK}_c(\mathbf{z})$ is $L$-positive:
\begin{equation}\label{eq:L-expansion}
  \widetilde{\GK}_c(\mathbf{z})
  = \sum_{\alpha} b_\alpha(c)\, L_\alpha,
  \qquad b_\alpha(c) \in \Z_{\geq 0}.
\end{equation}
\end{observation}

Since $M$-positivity is automatic (the $M$-coefficients count
CPPs), $L$-positivity is a strictly stronger assertion.
The $M$-to-$L$ conversion uses M\"obius inversion on the
refinement poset:
\[
  b_\alpha = \sum_{\beta \succeq \alpha}
  (-1)^{\ell(\beta) - \ell(\alpha)}\, a_\beta,
\]
where $\beta \succeq \alpha$ means $\alpha$ is a refinement of~$\beta$
(equivalently, $\beta$ is a coarsening of~$\alpha$).
The nonnegativity of the~$b_\alpha$ after this alternating sum is
nontrivial and requires cancellation.

\subsection{$L$-coefficients for $c = (1,1,0)$, $d = 2$}

\begin{center}
\begin{tabular}{c|l}
  \textbf{Composition $\alpha$} & \textbf{$b_\alpha$} \\
  \hline
  $()$ & $1$ \\
  $(k)$, any $k \geq 1$ & $2$ \\
  $(1, k)$, any $k \geq 1$ & $1$ \\
  $(j, k)$, $j \geq 2$ & $2$ \\
  $(1, 2, 1)$ & $1$ \\
  $(1, j, k)$, $j \geq 2$, $k \geq 2$ & $1$ \\
  $(j, 2, k)$, $j \geq 2$ & $2$ \\
  $(j, k, l)$, $j,k \geq 2$, $l \geq 2$, $k \geq 3$ & $2$
\end{tabular}
\end{center}

The $L$-coefficients are smaller than the $M$-coefficients and
exhibit even cleaner stabilization.

\subsection{Combinatorial interpretation}

$L$-positivity has a well-known combinatorial meaning:
a quasi-symmetric function is $L$-positive if and only if it
can be written as $\sum_w L_{D(w)}$, a sum of fundamental
quasi-symmetric functions indexed by descent sets, over some
set of ``words'' or ``permutations.''

Equivalently, by the work of Tewari and van
Willigenburg~\cite{tewari-vanwilligenburg}, $L$-positivity
of a quasi-symmetric function~$f$ is equivalent to~$f$
being the quasi-symmetric characteristic of a module
over the $0$-Hecke algebra $H_n(0)$.
Thus Observation~\ref{obs:L-pos} implies the existence
of a $0$-Hecke module whose character
is~$\widetilde{\GK}_c$.

\begin{remark}[The $0$-Hecke algebra]\label{rem:0hecke}
The \emph{$0$-Hecke algebra} $H_n(0)$ is the quotient
of the free algebra on generators $T_1, \ldots, T_{n-1}$
by the Coxeter braid relations together with the
\emph{idempotent} relations $T_i^2 = T_i$
(in contrast to the symmetric group algebra, where
$T_i^2 = 1$, or the Hecke algebra $H_n(q)$, where
$T_i^2 = (q-1)T_i + q$). Equivalently, $H_n(0)$
is obtained by specializing the Hecke algebra
$H_n(q)$ at $q = 0$.

The irreducible $H_n(0)$-modules are one-dimensional
and are indexed by compositions of~$n$ (not by
partitions, as for $\mathfrak{S}_n$). The
quasi-symmetric characteristic
$\operatorname{ch}(M) = \sum L_\alpha$, summed over
composition indices of the irreducible constituents,
provides the natural bridge between $H_n(0)$-modules
and $L$-positive quasi-symmetric functions.
\end{remark}

A related phenomenon appears in the work of Shareshian and
Wachs~\cite{shareshian-wachs}, who define a quasi-symmetric
refinement of the chromatic symmetric function and prove its
$L$-positivity (and even $e$-positivity) for certain classes
of graphs. Their construction is motivated by the
Stanley--Stembridge conjecture and is analogous in spirit:
a combinatorial generating function, refined by a
``descent-like'' statistic, turns out to be $L$-positive.

\begin{conjecture}\label{conj:descent}
There exists a natural descent-type statistic on cylindric
partitions such that $\widetilde{\GK}_c = \sum_\pi L_{\alpha(\pi)}$,
where $\alpha(\pi)$ is the descent composition of~$\pi$.
\end{conjecture}

Finding such a statistic would give a combinatorial proof of
$L$-positivity and could shed light on the original positivity
conjecture for~$Q_{n,c}$.

\subsection{Failure of the E--H factorization in QSym}

The Euler--Hamburger factorization $\GK_c = H \cdot F_c$,
where $H = 1/\qpoch{zq}{q}$ and $F_c = \qpoch{zq}{q} \cdot \GK_c$,
is the key identity connecting $\GK_c$ to the polynomials~$Q_{n,c}$.

We tested two natural QSym analogs:
\begin{enumerate}
  \item \textbf{With max-entry variable.} Introduce a variable $t$
    tracking the maximum entry and define
    $\widetilde{\GK}_c(t, \mathbf{z}) = \sum_n t^n \tilde{g}_n$,
    then compute $[t^n]\bigl(\prod_{k \geq 1}(1 - tz_k) \cdot
    \widetilde{\GK}_c(t, \mathbf{z})\bigr)
    = \sum_{j=0}^n (-1)^j e_j \cdot \tilde{g}_{n-j}$.
    This is \emph{not} quasi-symmetric and has \emph{negative}
    coefficients, because the individual
    $\tilde{g}_n$ (CPPs with max entry~$n$) are not quasi-symmetric:
    fixing the max entry privileges a particular value.

  \item \textbf{Without max-entry variable.} Compute
    $\prod_{k \geq 1}(1 - z_k) \cdot \widetilde{\GK}_c(\mathbf{z})$
    directly in $\QSym$, which under $z_k \mapsto q^k$
    specializes to $\qpochn{q}{q}{\infty} \cdot \GK_c(1,q)$.
    This also has \emph{negative} coefficients
    and is not quasi-symmetric.
\end{enumerate}

Both failures indicate that the E--H factorization is
fundamentally tied to the max-entry variable~$z$, which does
not sit naturally in the QSym framework.

%% ========================================================================
\section{Attempt to lift $Q_n$ to QSym: negative coefficients}
\label{sec:Qn-lift}
%% ========================================================================

The positive results above concern the \emph{full} generating function
$\widetilde{\GK}_c$, which sums over all cylindric partitions without
constraint on the maximum entry. The central positivity problem,
however, is about the \emph{polynomials} $Q_{n,c}(q)$, which are
extracted by the alternating-sign formula
\begin{equation}\label{eq:Q-ground-set}
  Q_{n,c}(q) = \sum_{j=0}^{n}\;
  \sum_{\substack{S \subseteq \{1,\ldots,n\} \\ |S|=j}}\;
  \sum_{U \subseteq \{1,\ldots,n-j\}}\;
  \sum_{\substack{\nu \in \CPP(c) \\ \max(\nu) = n-j}}
  (-1)^{|S|+|U|}\, q^{\sigma(S)+\sigma(U)+|\nu|}.
\end{equation}
The natural question is: can the QSym machinery, which works so well for
$\widetilde{\GK}_c$, be applied to $Q_{n,c}$?

\subsection{The QSym $Q$-polynomial}

Given a triple $(S, U, \nu)$ from the ground set
$\widetilde{\mathcal{G}}_{n,c}$, write $m = n - |S|$ and define the
\emph{combined frequency vector}:
\[
  \operatorname{freq}(S, U, \nu)_k
  = |\{s \in S : s = k\}|
  + |\{u \in U : u = k\}|
  + |\{(i,j) : \nu^{(i)}_j = k\}|,
\]
counting how many times the value~$k$ appears across all three
components. The total weight is
$\sigma(S) + \sigma(U) + |\nu| = \sum_k k \cdot \operatorname{freq}(S,U,\nu)_k$,
so replacing $q^k \mapsto z_k$ in each factor gives a
quasi-symmetric refinement:
\begin{equation}\label{eq:Q-qsym}
  Q_n^{\QSym}
  = \sum_{(S,U,\nu) \in \widetilde{\mathcal{G}}_{n,c}}
    (-1)^{|S|+|U|}\; M_{\operatorname{freq}(S,U,\nu)}.
\end{equation}
Under the principal specialization $z_k \mapsto q^k$,
$M_\alpha \mapsto \sum_{i_1 < \cdots < i_\ell} q^{\sum a_s i_s}$,
and summing over all frequency vectors recovers $Q_{n,c}(q)$.

\subsection{Failure: negative $M$-coefficients}

\begin{observation}[Computational]\label{obs:Q-not-M-pos}
For $n \geq 2$, the QSym polynomial $Q_n^{\QSym}$ has
\emph{negative} coefficients in the monomial basis~$M_\alpha$.
This holds for every tested profile with $d \leq 5$, including
the simplest case $c = (1,0,0)$, $d = 1$.
\end{observation}

Concretely, for $c = (1,1,0)$, $n = 2$:
\[
  Q_2^{\QSym}
  = M_{(1,1)} - 2\,M_{(2)} + M_{(1,1,1)} - M_{(2,1)} + \cdots
\]
The coefficient $a_{(2)} = -2$ is negative. Since $M$-positivity
fails, $L$-positivity fails \emph{a fortiori}: the $L$-basis
expansion also contains negative terms.

\begin{remark}
For $n = 1$, $Q_1^{\QSym}$ is $M$-positive and in fact equal to
$h_1^{\QSym}$, where $h_m$ is the ``corrected'' generating function
$h_m = \qpochn{q}{q}{m} \cdot g_m$ (cf.\ Section~\ref{sec:L-positivity}).
The failure begins precisely at $n = 2$, where the alternating
$q$-binomial sum introduces cancellation between different
frequency vectors.
\end{remark}

\subsection{Anatomy of the failure}

The key issue is \emph{cross-cancellation}: different frequency
vectors can map to the same $q$-weight under specialization,
so cancellation in $Q_n(q)$ \emph{must} cross QSym monomials.

More precisely, for a given $q$-weight~$w$, define:
\begin{itemize}
  \item The \emph{internal cancellation} at frequency $\alpha$:
    positive and negative triples with $\operatorname{freq} = \alpha$
    and total weight~$w$ that cancel within the same QSym class.
  \item The \emph{cross-cancellation}: excess positive contributions
    from one frequency $\alpha$ cancelling excess negative from
    another frequency $\beta$, where both have $|\alpha| = |\beta| = w$
    (same total weight, different composition).
\end{itemize}

Computationally, approximately 94--97\% of cancellation is internal
(same frequency vector, handled by the $U$-toggle involution).
But the remaining 0.3--3\% is cross-cancellation, which is
\emph{essential}: the net coefficient at some individual frequency
vectors is negative, and only the aggregation across all
compositions of the same weight yields the positive $Q_n$ coefficient.

\begin{example}
For $c = (1,1,0)$, $n = 2$, at $q$-weight~$4$:
the frequency vector $(0,2)$ (two copies of value~2)
has net coefficient $-1$, while $(1,0,1)$
(one copy of value~1 and one of value~3) has net $+2$.
Under $z_k \mapsto q^k$, both map to $q^4$,
and their contributions $-1 + 2 = 1$ yield $[q^4]\,Q_2 = 1 \geq 0$.
\end{example}

\subsection{The ``$\nu$-only'' variant}

One might hope that tracking only the CPP parts~$\nu$ in the
frequency vector (keeping $S$ and $U$ as scalar $q$-factors)
would avoid the problem:
\begin{equation}\label{eq:Q-nu-only}
  Q_n^{\nu\text{-only}}
  = \sum_{(S,U,\nu)}
    (-1)^{|S|+|U|}\, q^{\sigma(S)+\sigma(U)}
    \cdot M_{\operatorname{freq}(\nu)}.
\end{equation}
This is a hybrid object: a formal power series in~$q$ with
coefficients in $\QSym$.

\begin{observation}[Computational]\label{obs:nu-only-fails}
The $\nu$-only variant~\eqref{eq:Q-nu-only} also has negative
coefficients, and in fact fails \emph{worse} than the combined
version~\eqref{eq:Q-qsym}: it is negative already at $n = 1$,
for every profile tested.
\end{observation}

The failure at $n = 1$ is structural. When $n = 1$, the ground
set includes $(S = \{1\}, U = \emptyset, \nu = \emptyset)$
with sign $-1$ and $q$-scalar $q^1$, contributing
$-q \cdot M_{()}$ to $Q_1^{\nu\text{-only}}$. Meanwhile the
positive contributions from $(S = \emptyset, U, \nu)$ are spread
across many frequency vectors $M_\alpha$ with $\alpha \neq ()$.
There is no way for these to cancel the $-q \cdot M_{()}$ term,
since it lives in a different QSym component.

By contrast, the combined version~\eqref{eq:Q-qsym} works at
$n = 1$ precisely because $S = \{1\}$ contributes to
$M_{(1)}$ (the same QSym monomial as a CPP part equal to~1),
enabling internal cancellation.

\subsection{Diagnosis}

The QSym framework succeeds brilliantly for the \emph{unconstrained}
generating function $\widetilde{\GK}_c$, where every coefficient
is manifestly nonneg (it counts CPPs). The difficulty arises when
we try to \emph{extract} $Q_n$, which requires:
\begin{enumerate}
  \item The max-entry constraint $\max(\nu) = m$ (privileges a
    particular value, breaking the quasi-symmetric relabelling argument);
  \item The alternating sign $(-1)^{|S|+|U|}$ (introduces cancellation
    that crosses between different frequency vectors).
\end{enumerate}

The $q$-specialization $z_k \mapsto q^k$ ``forgets'' the
composition structure and allows cross-frequency cancellation to
produce the positive polynomial $Q_n(q)$. This forgetting is
essential and cannot be avoided within $\QSym$.

In summary: $\widetilde{\GK}_c \in \QSym$ with $L$-positivity
is a genuine structural property of cylindric partitions, but
it does not extend to the $Q$-polynomials. The positivity of
$Q_{n,c}(q)$ appears to be an \emph{intrinsically one-variable}
phenomenon.

%% ========================================================================
\section{Computational verification}
\label{sec:computation}
%% ========================================================================

The quasi-symmetry was verified as follows. For each profile~$c$
with $d \leq 5$:

\begin{enumerate}
  \item Enumerate all CPPs with maximum entry $\leq n_{\max}$ and
    total weight $\leq w_{\max}$, using the shifted dominance
    conditions~\eqref{eq:interlacing}.
  \item For each CPP, compute the frequency vector
    $(\#_1, \#_2, \ldots, \#_{n_{\max}})$ and extract the
    composition $\alpha$ (ordered nonzero entries).
  \item Group CPP counts by composition and verify that
    all placements of the same composition at different
    value positions give the same count (restricting to
    compositions whose weight does not exceed $w_{\max}$ at
    any placement, to avoid truncation artifacts).
\end{enumerate}

\begin{center}
\begin{tabular}{c|c|c|c}
  $d = |c|$ & Profiles tested & QSym failures & $L$-negative terms \\
  \hline
  1 & 3 & 0 & 0 \\
  2 & 6 & 0 & 0 \\
  3 & 10 & 0 & 0 \\
  4 & 15 & 0 & 0 \\
  5 & 21 & 0 & 0 \\
  \hline
  \textbf{Total} & \textbf{55} & \textbf{0} & \textbf{0}
\end{tabular}
\end{center}

Parameters used: $n_{\max} = 3$, $w_{\max} = 12$,
with 24 non-truncated quasi-symmetry checks per profile.
The $L$-basis conversion uses M\"obius inversion on the
refinement poset of compositions.

\medskip

The partial quasi-symmetry of
$\tilde{g}_n$ (Observation~\ref{obs:partial-qsym})
was verified separately by fixing the max entry~$n$
and checking that each $z_n^m$ slice is quasi-symmetric
in $z_1, \ldots, z_{n-1}$:

\begin{center}
\begin{tabular}{c|c|c|c}
  $d = |c|$ & Profiles tested & Partial QSym failures
  & Slice $L$-positive \\
  \hline
  1 & 3 & 0 & No \\
  2 & 6 & 0 & No \\
  3 & 10 & 0 & No \\
  4 & 15 & 0 & No \\
  5 & 21 & 0 & No \\
  \hline
  \textbf{Total} & \textbf{55} & \textbf{0} & \textbf{---}
\end{tabular}
\end{center}

Parameters: $n = 1, 2, 3$, $w_{\max} = 18$.
Partial quasi-symmetry holds for all 165 test cases
(55 profiles $\times$ 3 values of~$n$).
Slice $L$-positivity fails at $n = 3$ for every profile
(the slices $f_{n,m}$ have $M$-coefficients that decrease
with composition length, forcing negative $L$-coefficients).

%% ========================================================================
%% References
%% ========================================================================

\begin{thebibliography}{99}

\bibitem{gessel-multipartite}
I.~M.~Gessel,
Multipartite $P$-partitions and inner products of skew Schur functions,
in: \emph{Combinatorics and algebra (Boulder, Colo., 1983)},
Contemp.\ Math.\ \textbf{34}, Amer.\ Math.\ Soc., 1984, pp.~289--317.

\bibitem{shareshian-wachs}
J.~Shareshian and M.~L.~Wachs,
Chromatic quasisymmetric functions,
\emph{Adv.\ Math.}\ \textbf{295} (2016), 497--551.
arXiv:1405.4629.

\bibitem{stanley-ordered}
R.~P.~Stanley,
Ordered structures and partitions,
\emph{Mem.\ Amer.\ Math.\ Soc.}\ \textbf{119}, 1972.

\bibitem{tewari-vanwilligenburg}
V.~V.~Tewari and S.~J.~van Willigenburg,
Modules of the $0$-Hecke algebra and quasisymmetric Schur functions,
\emph{Adv.\ Math.}\ \textbf{285} (2015), 1025--1065.
arXiv:1403.1527.

\end{thebibliography}

\end{document}
